% ============================================================
%  Objetivo 1.1 — Comparación de 2 secuencias
% ============================================================
\section{Objetivo 1.1 — Comparación de 2 secuencias}

\textbf{¿Qué opciones de programas hay y para qué tipo de comparaciones sirven?}

En el formulario de BLAST aparecen cinco opciones:

\begin{itemize}
  \item \textbf{blastn}: busca una secuencia nucleotídica en una base de nucleótidos.
  \item \textbf{blastp}: busca una proteína en una base de proteínas (la que usamos en este ejercicio).
  \item \textbf{blastx}: traduce el query nucleotídico y lo busca en una base de proteínas.
  \item \textbf{tblastn}: busca una proteína en una base de nucleótidos que el programa traduce.
  \item \textbf{tblastx}: traduce tanto el query como la base y compara proteínas.
\end{itemize}

\textbf{¿Qué otras matrices hay disponibles? ¿Conoce sus diferencias?}

Además de BLOSUM62, NCBI BLAST ofrece PAM30, PAM70, PAM250, BLOSUM45, BLOSUM50, BLOSUM80 y BLOSUM90.
Las matrices BLOSUM se construyen a partir de bloques de secuencias alineadas: el número indica el porcentaje
de identidad promedio de las secuencias usadas para construirla. BLOSUM62 es la más general y la que mejor
funciona para la mayoría de las búsquedas. BLOSUM80 sirve cuando las secuencias a comparar son muy similares
entre sí (>80\% de identidad), porque puntúa mejor las sustituciones conservativas y penaliza más los cambios
drásticos. BLOSUM45 es la opción para secuencias más divergentes. Las matrices PAM son una familia distinta,
basada en modelos de evolución de aminoácidos, y se usan menos en la práctica.

\textbf{¿Qué particularidad tienen los gap costs? ¿A qué se debe?}

Los gap costs tienen dos componentes separados: \textbf{existence} (abrir el gap) y \textbf{extension}
(cada residuo adicional dentro del gap). Esto se conoce como modelo de gap afín y existe porque biológicamente
no es lo mismo abrir una inserción que extenderla. Abrir un gap corresponde a un evento evolutivo único
(un indel), mientras que agregar un residuo más dentro de ese mismo gap es menos costoso desde el punto de
vista evolutivo. Por eso la penalidad de apertura es mucho mayor que la de extensión (11 vs 1 en el default
de BLOSUM62). Además, los valores disponibles cambian según la matriz porque cada combinación tiene sus
propios parámetros estadísticos.

% ------------------------------------------------------------
\subsection{Ejercicio 1.1}

Comparamos la mioglobina humana (P02144, 154\,aa) con la de caballo (P68082, 154\,aa) y cerdo
(P02189, 154\,aa). Para todas las corridas usamos \textbf{blastp} desde el formulario
\textit{Align two or more sequences}.

\textbf{i) ¿El alineamiento abarca la secuencia completa de las 2 proteínas? ¿Le parece que es un
alineamiento local o global? Justifique.}

Sí, en los dos casos el alineamiento abarcó la secuencia completa (cobertura 100\%). BLAST es un algoritmo
de alineamiento local, pero cuando dos proteínas son tan similares como estas mioglobinas de mamíferos
($\sim$88\% de identidad), el único HSP significativo termina extendiéndose por toda la longitud de la
secuencia porque no hay ninguna región de baja similitud que haga conveniente cortarlo antes. El resultado
se ve prácticamente igual a un alineamiento global, aunque algorítmicamente sigue siendo local.

\textbf{ii) ¿Se observan gaps en el dot-plot? ¿Cuántos? ¿Cómo se identifican?}

No hay gaps en ninguna de las dos comparaciones. El dot-plot muestra una única línea diagonal recta y
continua de la posición 1 a la 154, sin ningún quiebre ni desplazamiento lateral, lo que es consistente
con el 0/154 (0\%) de gaps del alineamiento. Si hubiera gaps, la diagonal mostraría saltos o segmentos
desplazados en paralelo: un gap en la query desplazaría la línea hacia abajo, un gap en el hit hacia la
derecha. En este caso la diagonal es perfecta.

\textbf{iii) Analice los valores de Score, E-value, Identities, Positives y Gaps.
¿Qué cambia al modificar la matriz y los gap costs?}

Con parámetros default (BLOSUM62, gap costs 11:1):

\begin{table}[H]
\centering
\footnotesize
\begin{tabular}{lP{2.2cm}P{2.4cm}P{2.4cm}P{2.4cm}P{1.8cm}}
\toprule
Organismo & Score bits (raw) & E-value & Identities & Positives & Gaps \\
\midrule
Caballo (P68082) & 278 (711) & $7 \times 10^{-103}$ & 136/154 (88\%) & 142/154 (92\%) & 0/154 (0\%) \\
Cerdo (P02189)   & 276 (705) & $6 \times 10^{-102}$ & 136/154 (88\%) & 144/154 (93\%) & 0/154 (0\%) \\
\bottomrule
\end{tabular}
\caption{BLASTp mioglobina humana vs.\ caballo y cerdo (BLOSUM62, gap costs 11:1).}
\end{table}

Los dos alineamientos muestran una homología muy alta: E-values cercanos a 0, identidad del 88\% y ningún
gap. Los Positives (92--93\%) superan las Identities porque varias de las posiciones que difieren son
sustituciones conservativas, es decir, cambios entre aminoácidos de propiedades similares. Algo que nos
llamó la atención: caballo y cerdo tienen exactamente la misma identidad (88\%) pero el cerdo tiene un
Positive más (93\% vs 92\%), lo que indica que sus sustituciones son ligeramente más conservativas que las
del caballo.

Para analizar el efecto de los parámetros usamos humana vs.\ caballo como referencia:

\begin{table}[H]
\centering
\footnotesize
\begin{tabular}{lP{2.2cm}P{2.4cm}P{2.4cm}P{2.4cm}P{1.8cm}}
\toprule
Configuración & Score bits (raw) & E-value & Identities & Positives & Gaps \\
\midrule
BLOSUM62, 11:1 (default) & 278 (711) & $7 \times 10^{-103}$ & 136/154 (88\%) & 142/154 (92\%) & 0/154 (0\%) \\
BLOSUM80, 10:1           & 291 (666) & $3 \times 10^{-102}$ & 136/154 (88\%) & 144/154 (93\%) & 0/154 (0\%) \\
BLOSUM62, 11:2           & 308 (711) & $5 \times 10^{-111}$ & 136/154 (88\%) & 142/154 (92\%) & 0/154 (0\%) \\
BLOSUM62, 9:1            & 217 (711) & $3 \times 10^{-84}$  & 136/154 (88\%) & 142/154 (92\%) & 0/154 (0\%) \\
\bottomrule
\end{tabular}
\caption{Efecto de la matriz y los gap costs en el alineamiento humana vs.\ caballo.}
\end{table}

\textbf{a) BLOSUM80:} el raw score bajó de 711 a 666 porque BLOSUM80 tiene valores numéricos distintos
para cada par de aminoácidos. Pero el score en bits subió (291 vs 278) y los Positives aumentaron
levemente (93\% vs 92\%). Esto tiene sentido porque BLOSUM80 fue construida con secuencias más similares
entre sí, entonces es más sensible a las sustituciones conservativas. El E-value empeoró un poco
($3 \times 10^{-102}$ vs $7 \times 10^{-103}$) por diferencias en los parámetros estadísticos
$\lambda$ y $K$ propios de esta matriz.

\textbf{b) Gap costs 11:2 (mayor penalidad de extensión):} el raw score se mantuvo en 711 porque el
alineamiento ya tenía 0 gaps, entonces penalizar más la extensión no cambia nada en el alineamiento en sí.
Lo llamativo es que el score en bits subió bastante (308) y el E-value mejoró mucho
($5 \times 10^{-111}$ vs $7 \times 10^{-103}$). Esto pasa porque con gaps más costosos, la probabilidad
de encontrar un alineamiento igual de bueno por azar es menor, haciendo el resultado estadísticamente
más significativo.

\textbf{c) Gap costs 9:1 (menor penalidad de extensión):} mismo razonamiento --- el raw score quedó en 711
porque el alineamiento no tiene gaps. Pero el score en bits cayó fuerte (217) y el E-value empeoró bastante
($3 \times 10^{-84}$ vs $7 \times 10^{-103}$). Al bajar el costo de los gaps, las secuencias aleatorias
pueden alcanzar puntajes similares más fácilmente, con lo cual el mismo resultado vale estadísticamente
mucho menos.

% ============================================================
%  Objetivo 1.2 — Búsqueda de proteínas homólogas en humanos
% ============================================================
\section{Objetivo 1.2 — Búsqueda de proteínas homólogas a la Mb en humanos}

\subsection{Ejercicio 1.2}

\textbf{i) ¿Cómo mejoraría la búsqueda para obtener solo las homólogas a la Mb una sola vez?}

Usando \textbf{RefSeq Select proteins} con el organismo restringido a \textit{Homo sapiens}. RefSeq Select
tiene una sola entrada representativa por gen, lo que elimina todas las entradas duplicadas que aparecen en
nr: distintas estructuras cristalográficas de la misma proteína, isoformas con accession diferente,
secuencias redundantes de distintos proyectos de secuenciación. Con nr sin filtro, la mioglobina sola puede
aparecer docenas de veces.

\textbf{ii) Describa brevemente cada proteína homóloga encontrada. Compare E-values, scores y
alineamientos.}

Con BLASTp contra RefSeq Select en \textit{Homo sapiens} con parámetros default encontramos 3 proteínas:

\begin{table}[H]
\centering
\footnotesize
\begin{tabular}{llP{2cm}P{2cm}P{1.5cm}P{1.5cm}}
\toprule
Proteína & Accesion & Score (bits) & E-value & Cobertura & Identidad \\
\midrule
myoglobin isoform 1          & NP\_005359.1 & 312  & $1 \times 10^{-111}$  & 100\% & 100\%   \\
cytoglobin (CYGB)            & NP\_599030.1 & 73.6 & $5 \times 10^{-17}$   & 95\%  & 28.67\% \\
hemoglobin subunit zeta (HBZ) & NP\_005323.1 & 71.6 & $1 \times 10^{-16}$   & 96\%  & 28.19\% \\
\bottomrule
\end{tabular}
\caption{Proteínas homólogas a la mioglobina humana (BLASTp, RefSeq Select, \textit{Homo sapiens}).}
\end{table}

La mioglobina misma aparece primero con 100\% de identidad (hit trivial, es la misma proteína). La
citoglobina y la hemoglobina zeta tienen identidades muy bajas ($\sim$28\%), lo que refleja la divergencia
evolutiva de las globinas a lo largo de cientos de millones de años. Aun así, los E-values son
significativos ($5 \times 10^{-17}$ y $1 \times 10^{-16}$), lo que confirma que hay homología real
aunque sea lejana. La diferencia de score entre la Mb (312 bits) y las otras dos (73--71 bits) es enorme,
y se ve también en los alineamientos: con la Mb el alineamiento es perfecto, mientras que con CYGB y HBZ
hay muchas diferencias y algunos gaps.

\textbf{iii) ¿Cuál es la globina humana más parecida a la Mioglobina? ¿Le sorprende algún resultado?}

La más parecida es la \textbf{citoglobina (CYGB)}, con E-value $5 \times 10^{-17}$ y 28.67\% de identidad.
Tiene sentido porque CYGB es, como la Mb, una globina intracelular monómera --- a diferencia de las
hemoglobinas, que circulan en la sangre como tetrámeros. Lo que nos sorprendió un poco es que la
hemoglobina zeta (HBZ) aparezca con una identidad tan similar a la de CYGB (28.19\%). HBZ es una
hemoglobina embrionaria muy antigua y divergente, con diferencias tan grandes respecto a las otras Hb
que termina pareciéndose a la Mb casi tanto como la citoglobina.

\textbf{iv) ¿Creen que encontraron todas las proteínas de la familia de las globinas que existen en el
genoma humano?}

No. Con BLASTp y RefSeq Select solo aparecieron 3 proteínas. Sabemos que faltan al menos la neuroglobina
(NGB), las hemoglobinas alfa (HBA1), beta (HBB), delta (HBD), épsilon (HBE1), gamma-1 (HBG1),
gamma-2 (HBG2), mu (HBM) y theta-1 (HBQ1). Algunas de estas probablemente quedan por debajo del umbral
de E-value con los parámetros por defecto dado lo baja que es su identidad con la Mb.

\textbf{v) Intente encontrar la mayor cantidad de globinas modificando los parámetros. ¿Qué parámetros
modificaron y qué proteínas nuevas encontraron?}

Con PSI-BLAST (ver punto vi) recuperamos prácticamente todas las globinas sin tocar los parámetros de
scoring. Si quisiéramos hacerlo con BLASTp simple, los cambios que probaríamos son: aumentar el E-value
threshold a 10 para no perder hits distantes, cambiar la matriz a BLOSUM45 (más adecuada para secuencias
con $\sim$45\% de identidad como estas globinas) y reducir el word size a 2 para que BLAST encuentre más
palabras semilla en secuencias divergentes.

\textbf{vi) PSI-BLAST: ¿encontraron alguna proteína nueva? ¿Cuántas iteraciones hicieron?}

Hicimos \textbf{2 iteraciones} de PSI-BLAST con la Mb humana como query contra RefSeq Select en
\textit{Homo sapiens}.

\medskip
\noindent\textbf{Iteración 1} --- 9 secuencias con $E < 0.005$:

\begin{table}[H]
\centering
\footnotesize
\begin{tabular}{lP{2.4cm}P{1.8cm}c}
\toprule
Proteína & E-value & Identidad & Nueva respecto a BLASTp \\
\midrule
myoglobin isoform 1              & $1 \times 10^{-111}$ & 100\%   & no \\
cytoglobin (CYGB)                & $5 \times 10^{-17}$  & 28.67\% & no \\
hemoglobin subunit zeta (HBZ)    & $1 \times 10^{-16}$  & 28.19\% & no \\
hemoglobin subunit theta-1 (HBQ1)& $5 \times 10^{-10}$  & 25.00\% & sí \\
hemoglobin subunit gamma-2 (HBG2)& $5 \times 10^{-10}$  & 25.55\% & sí \\
hemoglobin subunit gamma-1 (HBG1)& $7 \times 10^{-10}$  & 25.55\% & sí \\
hemoglobin subunit alpha (HBA1)  & $6 \times 10^{-9}$   & 27.52\% & sí \\
hemoglobin subunit epsilon (HBE1)& $4 \times 10^{-8}$   & 24.09\% & sí \\
hemoglobin subunit mu (HBM)      & $7 \times 10^{-8}$   & 27.61\% & sí \\
\bottomrule
\end{tabular}
\caption{PSI-BLAST iteración 1: globinas humanas encontradas.}
\end{table}

Ya en la primera iteración aparecieron 6 globinas nuevas que BLASTp no había encontrado.

\medskip
\noindent\textbf{Iteración 2} --- 12 secuencias en total, 3 nuevas:

\begin{table}[H]
\centering
\footnotesize
\begin{tabular}{lP{2.4cm}P{1.8cm}}
\toprule
Proteína nueva & E-value & Identidad \\
\midrule
hemoglobin subunit delta (HBD) & $6 \times 10^{-57}$ & 25.52\% \\
hemoglobin subunit beta (HBB)  & $9 \times 10^{-56}$ & 24.83\% \\
neuroglobin (NGB)              & $6 \times 10^{-14}$ & 17.76\% \\
\bottomrule
\end{tabular}
\caption{PSI-BLAST iteración 2: proteínas nuevas respecto a la iteración 1.}
\end{table}

El hallazgo más interesante es la \textbf{neuroglobina (NGB)} en la segunda iteración, con solo 17.76\%
de identidad respecto a la Mb. BLASTp nunca la hubiera encontrado con esos parámetros porque esa identidad
es demasiado baja para que el hit sea estadísticamente significativo usando BLOSUM62. PSI-BLAST pudo
recuperarla porque después de la primera iteración construye una \textbf{PSSM}
(\textit{Position-Specific Scoring Matrix}) con todas las globinas encontradas. La PSSM asigna un puntaje
diferente a cada aminoácido en cada posición de la secuencia según cuánto está conservada esa posición en
la familia --- es mucho más informativa que una matriz genérica como BLOSUM para detectar homólogos lejanos.
